% FDUletter_example.tex - an example latex file to illustrate FDUletter.cls
% OUCletter_example.tex - an example latex file to illustrate OUCletter.cls
%
% Template by Brian Wood (brian.wood@oregonstate.edu).  Please feel free to send suggestions for changes; this template/cls is not exactly elegantly done!
% Modified by Huang Weiran (huangweiran1998@outlook.com) to fit the need of FDU students.
% Modified by Tim (qt@stu.ouc.edu.cn) to fit the need of OUC students.


\documentclass[12pt]{OUCletter}
\usepackage{fontspec} 
\usepackage{tikz} 
\usepackage{xcolor}
	\definecolor{scured}{RGB}{1, 45, 114}
\usepackage{lipsum}
\usepackage{fancyhdr}
\usepackage{lastpage}
\usepackage{eso-pic}
%
% This section is just a bunch of busywork so that the second and following pages read ``Page X of Y''
\pagestyle{fancy}
\fancyhf{}
\renewcommand{\headrulewidth}{0pt}
\renewcommand{\footrulewidth}{0pt}
\rhead{Page \thepage \hspace{1pt} of \pageref{LastPage}}
%
%
% Set custom font here. Comment this line out if you do not have a Cambria font (originally included with this template) installed; computer modern (or whatever your current default font is) will be substituted.
%
%\setmainfont{[Cambria.ttf]}[BoldFont  = [CambriaBold.ttf], ItalicFont  = [CambriaItalic.ttf], BoldItalicFont = [CambriaBoldItalic.ttf] ]

% \newcommand{\watermark}[3]{\AddToShipoutPictureBG{
% \parbox[b][\paperheight]{\paperwidth}{
% \vfill%
% \centering%
% \begin{tikzpicture}
%     \path (0,0) -- (\paperwidth,\paperheight);
%     \node[opacity=.05] at (current page.center)
%     {\includegraphics[width=0.55\textwidth]{background-OUC.pdf}};   % sculogo watermark
%     \end{tikzpicture}
% \vfill}}}

% The material below is a whole big dang thing whose purpose is just to set up a fixed coordinate system for \tikz so that you can put the Department or School address in the upper right-hand side without it moving all around every time you change something in the page.  I think it works.
% Defining a new coordinate system for the page:
%
% --------------------------
% |(-1,1)    (0,1)    (1,1)|
% |                        |
% |(-1,0)    (0,0)    (1,0)|
% |                        |
% |(-1,-1)   (0,-1)  (1,-1)|
% --------------------------
\makeatletter
\def\parsecomma#1,#2\endparsecomma{\def\page@x{#1}\def\page@y{#2}}
\tikzdeclarecoordinatesystem{page}{
	\parsecomma#1\endparsecomma
	\pgfpointanchor{current page}{north east}
	% Save the upper right corner
	\pgf@xc=\pgf@x%
	\pgf@yc=\pgf@y%
	% save the lower left corner
	\pgfpointanchor{current page}{south west}
	\pgf@xb=\pgf@x%
	\pgf@yb=\pgf@y%
	% Transform to the correct placement
	\pgfmathparse{(\pgf@xc-\pgf@xb)/2.*\page@x+(\pgf@xc+\pgf@xb)/2.}
	\expandafter\pgf@x\expandafter=\pgfmathresult pt
	\pgfmathparse{(\pgf@yc-\pgf@yb)/2.*\page@y+(\pgf@yc+\pgf@yb)/2.}
	\expandafter\pgf@y\expandafter=\pgfmathresult pt
}
\makeatother


%%------------------------------------------------------------------------------%%
%

%%%%%%%%%%% Put Recommender(your mentor) Information Here %%%%%%%%%%%
%
\def\name{\textbf{Jie Gao} \\
	Malone Postdoctoral Fellow \\
	Department of Computer Science \\
    Malone Center for Engineering in Healthcare \\
    Center for Language and Speech Processing \\
    jgao77@jh.edu \\
    https://gaojie058.github.io/ \\
    
}
%
% List your degree(s), licences, etc. here if you like.
%\def\What{, Your degrees, etc.} 
%
% Set the name of your Department or School here
% I honestly don't know why the negative spacing is necessary to get the alignment of the first line correct.  This must be a ``\tikz'' thing.
%%%%%%%%%%%%%%%%%%  School or Department %%%%%%%%%%%%%%%
\def\Where{\hspace{-1.2mm}\textbf{\color{scured}%{scured}
		Department of Computer Science, \\
        Johns Hopkins University 
}} 

%%%%%%%%%%%%  Additional Contact Information %%%%%%%%%%%
%
% Set your preferred primary contact address here.
\def\Address{Malone Hall, 3400 N Charles St} 
%
\def\CityZip{Baltimore, MD 21218, USA %\\
%	P.R.China
} 
%
% Set your OUC e-mail here
\def\Email{\textbf{\color{scured}E-mail}: jgao77@jh.edu}
%
% Set your preferred contact number here
\def\TEL{\textbf{\color{scured}Phone}: +1 (410)710-0555}
%
% Set your department or personal website here
\def\URL{\textbf{\color{scured}Website}: {https://gaojie058.github.io/}}
%

%%%%%%%%%%%%%%%%  Footer information  %%%%%%%%%%%%%%%%%%
%
%  The next line is for your college, used as a footer.  If you prefer not to have this, just comment out these lines in favor of the line labeled "[[Alternate]]" below
\def\school{\small{
		OUC $\cdot$
		~Department of xxxxx $\cdot$
		~No. 238 Songling Road $\cdot$
		~Qingdao, China P.R} } 
% \def\school{~}  % [[Alternate]]
%

%%%%%%%%%%%%%%%%%%%%%  Signature line  %%%%%%%%%%%%%%%%%%%%%
%
% Set your signature line here.
% One can add a signature image in a PDF file using the following code; this requires a file called "signature_block.pdf" to be installed in the same folder as the .tex file.  The vertical spacing (\vspace) and the scaling will have to be adjusted to get things to look correct for your particular signature image. Alternatively, comment out the following line in favor of the one labeled "[[Alternate]]" if you want to sign a paper copy of the letter.
%
\signature{
	\vspace{2mm}\includegraphics[scale=0.8]{signature.pdf}\\\vspace{2mm}
	\name
}
%\signature{\name}  % [[Alternate]]
%%------------------------------------------------------------------------------%%



% This block sets up the address on the right-hand side of the header. 
%
% The following lines just compile the information you set up into the LaTex letter variable "address" for later use.
%
%The following command "clears out" the default address so that it can be better set using \tikz
\address{}

\def\newaddress{
	\Where\\ 
	\Address\\ 
	\CityZip\\ 
	\TEL\\ 
	\Email\\ 
	\URL 
}
%
%%%%%%%%%%%  DATE  %%%%%%%%%%%%%%%%%%%%%%%%%
% If you want a date different from the current date, comment out the next line in favor of the line labeled "[[Alternate]]".  The ``\vspace{10mm}'' just moves the date down a tiny bit so it doesn't interfere with the header.  This can be adjusted to your preference.
%
\date{\vspace{10mm} \today}
%\date{\vspace{10mm} 20 September 2020}  %[[Alternate]]
%
%%%%%%%%%%% Set the subject here if there is one  %%%%
%\subject{Stuff} % optional subject line



%%-----------------------------------------------------------------------------------------%%
%
\begin{document}
	%
	%
	%%%%%%%%  The "To" address goes here.
	%
	\begin{letter}{
			\textit{Faculty Search Committee} \\
            Digital Security, Safety, and Trust, School of Information \\
            University of California, Berkeley
		%	IOP Publishing%\\ 
			%SomeTown, SomeState 					       				  		 ~~SomeZip
		}
		% This line sets up the return address to the right-side of the OSU logo.  The location is set with absolute node addresses using ``\tikz''.  It can still be a bit fussy, and you may need to alter this a little to get things to look right.  The bit that changes the position are the numbers in parentheses ``at (14.2,2.7)''
		%
		\begin{tikzpicture}[remember picture,overlay,,every node/.style={anchor=center}]
		\node[text width=7cm] at (page cs:0.5,0.73){\small \newaddress};
		\end{tikzpicture} 
		
		%%%%%%  The ``opening'' is just the method of address you would like to use at the start of the letter.
		%
		\opening{Dear Members of Search Committee,}
		
		
		%%%%%%%%%% Body of letter   %%%%%%%%%%%%%%
		% Remove it if you do not want watermark
		% \watermark{}{}{}
		
		%%%%%%%%%%% 	Text   %%%%%%%%%%%%%%
		It is with great enthusiasm that I submit my application for the \textbf{Tenure-Track Assistant Professor position at University of California, Berkeley}. I am currently a postdoctoral researcher at Johns Hopkins University, where my research focuses on \textbf{human–AI collaboration}, \textbf{responsible and trustworthy AI}, and \textbf{computational social science}. My interdisciplinary background in human–computer interaction (HCI) and natural language processing (NLP) provides a unique human-centric perspective that closely aligns with the University of California, Berkeley’s vision for shaping and advancing the field of Security, Safety, and Trust in the context of human interaction with digital technologies.

        At Johns Hopkins University, I collaborate with researchers in HCI, NLP, and Robotics to investigate how \textit{trustworthiness}, \textit{transparency}, and \textit{human agency} can be implemented and facilitated in LLM-assisted workflows. With a focus on human and data, I study how AI can be used to empower humans to understand and utilize complex, unstructured data, such as interview transcripts, software code, and social media text, in a rigorous, transparent, and trustworthy way. My research seeks to answer a fundamental question: \textit{How can we design AI systems that empower, rather than replace, human capability that learns from data?} My long-term research vision is to expand human capabilities through human–AI collaboration, extending beyond the domain of data analysis into broader areas, enabling responsible and explainable AI use across domains such as science, healthcare, and policy. My projects, develop LLM-powered systems that bring humans into the AI loop, \textit{CollabCoder (CHI2024)} and \textit{CoAIcoder (TOCHI2023)}, making AI-assisted data analysis both efficient and trustworthy (MindCoder (2025). Two of them have been open source (CollabCoder) and publicly available now (mindcoder.ai). My research advances the foundations and practice of human-centered data science by integrating methods from NLP, software engineering (SE), and social science (SS). It also contributes to the growing literature on safe and transparent AI systems by operationalizing classical trustworthiness principles (credibility, dependability, transferability, and confirmability) into AI workflows.

        Thank you for considering my application. I am confident that my research can contribute meaningfully to UCB’s information community, and I would be honored to discuss how my research and teaching align with UCB’s vision for security, safety, and trust in the digital era. 

		
		%%%%%%% ``closing'' sets the sign-off line.
		\closing{Sincerely yours,}
		
		% Comment out/in the lines below as necessary
		%\encl{If an enclosure is provided, let them know what it is.}
		
		%\ps{A postscript if that is a thing you do.}
		
		%\cc{Someone Who Cares (and is copied).}
		
	\end{letter}
	
\end{document}
%
%%-----------------------------------------------------------------------------------------%%